% Section content template for individual Educative lessons
% This file contains the actual content for a specific section
% Generated from Educative API section content

% Section: Code of Library Management System
% Section ID: 6101410175516672
% Section Slug: code-of-library-management-system
% Generated: 2025-12-11T15:13:37.472402

We've gone over the different aspects of the library management system and observed the attributes attached to the problem using various UML diagrams. Let us now explore the more practical side of things, where we will work on implementing the library management system using multiple languages. This is usually the last step in an object-oriented design interview process.
\\
We have chosen the following languages to write the skeleton code of the different classes present in the library management system:

\begin{itemize}
\item
\protect\phantomsection\label{5IqdwXmRp5jvqrxHtM1qW}
Java
\item
\protect\phantomsection\label{FehgZxJfaHT18MN8X-_Je}
C\#
\item
\protect\phantomsection\label{NPb7JgvteNP2-FO-1jUct}
Python
\item
\protect\phantomsection\label{zEbwRTKlsYxax2BBeyYIq}
C++
\item
\protect\phantomsection\label{vgdFm6wABiusYgxaEx9bp}
JavaScript
\end{itemize}
\\
If you want to skip directly to the implementation and see the complete workflow, simply scroll down or click \hyperref[Executable-code-Library-management-system]{\textbf{Executable Code: Library Management System}}.

\subsection{Library management}\label{sTgw9YDUl0ePssZT9QiUR}
\\
This section will provide the skeleton code of the classes designed in the class diagram lesson.
\begin{quote}
\textbf{Note:} For simplicity, we are not defining getter and setter functions. The reader can assume that all class attributes are private, accessed through their respective public getter methods, and modified only through their public method functions.
\end{quote}
\subsubsection{Enumerations}\label{ahBFmnf5KAWv620-Swp37}
\\
First, we will define all the enumerations required in the library management system. According to the class diagram, four enumerations are used in the system: \texttt{BookFormat}, \texttt{BookStatus}, \texttt{ReservationStatus}, and \texttt{AccountStatus}. The code to implement these enumerations is as follows:`
\begin{quote}
\textbf{Note:} JavaScript does not support enumerations, so we will use the \texttt{Object.freeze()} method as an alternative that freezes an object and prevents further modifications.
\end{quote}

\textbf{Enum definitions}

\begin{lstlisting}[language=Java]
// definition of enumerations used in library management system enum BookFormat { HARDCOVER,
PAPERBACK, AUDIOBOOK,
EBOOK, NEWSPAPER,
MAGAZINE, JOURNAL
}

enum BookStatus { AVAILABLE,
RESERVED, LOANED,
LOST }

enum ReservationStatus { WAITING,
PENDING, CANCELED,
NONE }

enum AccountStatus { ACTIVE,
CLOSED, CANCELED,
BLOCKLISTED, NONE
}
\end{lstlisting}

\subsubsection{Address and person}\label{4zRC4VVx2apOCZH6RDWtH}
\\
This section contains the code for \texttt{Address} and \texttt{Person} classes, where the \texttt{Person} class is composed of an \texttt{Address} class. The implementation of these classes can be found below:

\textbf{The Address and Person classes}

\begin{lstlisting}[language=Java]
public class Address { private String streetAddress;
private String city; private String state;
private int zipCode; private String country;
}

public class Person { private String name;
private Address address; private String email;
private String phone; }
\end{lstlisting}

\subsubsection{User}\label{HpiNtZKLfaAvqo7GfYfmi}
\\
The \texttt{User} is an abstract class that represents the various people or actors that can interact with the system. As there are two types of users, the librarian and the library member, the user can either be a \texttt{Librarian} or a \texttt{Member}. The implementation of the mentioned classes is shown below:

\textbf{User and its child classes}

\begin{lstlisting}[language=Java]
// User is an abstract class public abstract class User { private String id;
private String password; private AccountStatus status;
private Person person; private LibraryCard card;

public abstract boolean resetPassword(); }

public class Librarian extends User { public boolean addBookItem(BookItem bookItem);
public boolean blockMember(Member member); public boolean unBlockMember(Member member);
public boolean resetPassword() { // definition
 }
public boolean addBookItem(BookItem bookItem) { // definition
 }
}

public class Member extends User { private Date dateOfMembership;
 private int totalBooksCheckedOut;

public boolean reserveBookItem(BookItem bookItem); private void incrementTotalBooksCheckedOut();
public boolean checkoutBookItem(BookItem bookItem); private void checkForFine(String bookItemId);
public void returnBookItem(BookItem bookItem); public boolean renewBookItem(BookItem bookItem);
public boolean resetPassword() { // definition
 }
}
\end{lstlisting}

\subsubsection{Book reservation, book lending, and fine}\label{yFzjoJZWw97_GO0EvJiWG}
\\
This component shows the implementation of \texttt{BookReservation}, \texttt{BookLending}, and Fine\ classes. These classes will be responsible for managing reservations against books, managing reservations, and calculating fines on books. The code is shown below:

\textbf{The BookReservation, BookLending, and Fine classes}

\begin{lstlisting}[language=Java]
public class BookReservation { private String itemId;
private Date creationDate; private ReservationStatus status;
 private String memberId;

public static BookReservation fetchReservationDetails(String bookItemId); }

public class BookLending { private String itemId;
private Date creationDate; private Date dueDate;
private Date returnDate; private String memberId;
private BookReservation bookReservation; private User user;

public static boolean lendBook(String bookItemId, String memberId); public static BookLending fetchLendingDetails(String bookItemId);
}

public class Fine { private Date creationDate;
private String bookItemId; private String memberId;

public static void collectFine(String memberId, long days); }
\end{lstlisting}

\subsubsection{Book and rack}\label{D_SSrTz9SQkmCH7Xg5sRO}
\\
The \texttt{Book} is an abstract class, and \texttt{BookItem} represents each book copy. For example, if there are two copies of the same book, there would be only one \texttt{Book} object and two \texttt{BookItem} objects. The code to implement these classes is as follows:

\textbf{The Book, BookItem, and Rack classes}

\begin{lstlisting}[language=Java]
// User is an abstract class
public abstract class Book {
    private String isbn;
    private String title;
    private String subject;
    private String publisher;
    private String language;
    private int numberOfPages;
    private BookFormat bookFormat;
    private List<Author> authors;
}

public class BookItem {
    private String id;
    private boolean isReferenceOnly;
    private Date borrowed;
    private Date dueDate;
    private double price;
    private BookStatus status;
    private Date dateOfPurchase;
    private Date publicationDate;
    private Rack placedAt;
    private Book book;  // Agggregation: BookItem has a reference to a Book
    
    // Constructors, getters, and other existing methods...

    public boolean checkout(String memberId) {
        // Implementation for checkout logic
        // Update the status, set due date, etc.
        // Return true if checkout is successful, false otherwise
        return true;  // Placeholder, replace with actual logic
    }

    public void setPlacedAt(Rack rack) {
        this.placedAt = rack;
        // Additional logic if needed
    }

    public void setAddedBy(Librarian librarian) {
        // Implementation for setting the librarian who added the book item
        // This might involve updating logs or other data related to the librarian
        // No return value as it's a setter method
    }
    
    // Other methods...
    
    public Book getBook() {
        return book;
    }
}

public class Rack {
    private int number;
    private String locationIdentifier;
    private List<BookItem> bookItems;
    public void addBookItem(BookItem bookItem) {
        // definition
    }
}

\end{lstlisting}

\subsubsection{Notification}\label{hS1Z6of4z6gsN_qmxF_4R}
\\
The \texttt{Notification} class is another abstract class responsible for sending notifications to the users, with the \texttt{PostalNotification} and \texttt{EmailNotification} classes as its child classes. The implementation of this class can be found below:

\textbf{Notification and its derived classes}

\begin{lstlisting}[language=Java]
// User is an abstract class public abstract class Notification { Private String notificationId;
Private Date creationDate; Private String content;
private BookLending bookLending; private BookReservation bookReservation;

public abstract boolean sendNotification(); }

public class PostalNotification extends Notification { private Address address;

public boolean sendNotification(){ // definition
 }
}

public class EmailNotification extends Notification { private String email;

public boolean sendNotification(){ // definition
 }
}
\end{lstlisting}

\subsubsection{Search and catalog}\label{i0qs6vex-wnCFWs7rSxek}
\\
The \texttt{Search} is an interface used to efficiently search library books by various methods, and the \texttt{Catalog} class is used to implement the search interface to help in book searching. The code to perform this functionality is presented below:

\textbf{The Search interface and the Catalog class}

\begin{lstlisting}[language=Java]
public interface Search { // Interface method (does not have a body)
public List<Book> searchByTitle(String title); public List<Book> searchByAuthor(String author);
public List<Book> searchBySubject(String subject); public List<Book> searchByPublicationDate(Date publishDate);
}

public class Catalog implements Search { private HashMap<String, List<Book>> bookTitles;
private HashMap<String, List<Book>> bookAuthors; private HashMap<String, List<Book>> bookSubjects;
 private HashMap<String, List<Book>> bookPublicationDates;

public List<Book> searchByTitle(String query) { // definition
 }
public List<Book> searchByAuthor(String query) { // definition
 }
public List<Book> searchBySubject(String query) { // definition
 }
public List<Book> searchByPublicationDate(String query) { // definition
 }
}
\end{lstlisting}

\subsubsection{Library}\label{bJZ4KI7Xvxf8FotChFLrZ}
\\
The final class of LMS is the \texttt{Library} class, which will be a Singleton class, meaning the entire system will have only one instance of this class. The implementation of this class can be found below:

\textbf{The Library class}

\begin{lstlisting}[language=Java]
public class Library {
private String name; private Address address;

// Private constructor to prevent external instantiation private Library() {}

 public Address getAddress();

// The Library is a singleton class that ensures it will have only one active instance at a time private static Library library = null;

// Created a static method to access the singleton instance of Library class public static Library getInstance () {
if (library == null) {
library = new Library (); }
return library; }
}
\end{lstlisting}

\subsection{Executable code: Library management system}\label{AQh0LIlifDDzcwQ4hpIzg}
\\
Below is a fully self-contained, runnable program in Java, C\#, C++, Python, and JavaScript that demonstrates the core workflows of the library management system. The system's main driver code resides in \texttt{Driver.java}, \texttt{Driver.cs}, \texttt{Driver.py}, \texttt{Driver.cpp}, and \texttt{Driver.js} for all respective languages. You can click the ``Run'' button to execute the code.

\subsubsection{What does this code show}\label{04XkVISpCfhbG-6N8yGlK}

\begin{itemize}
\item
\protect\phantomsection\label{7EFXP5upV1bUsIv9mBy_E}
\textbf{System initialization:} Sets up books, copies, members, a librarian, and the catalog.
\item
\protect\phantomsection\label{lapO-FuT5TaGSEZ3qxHHU}
\textbf{Scenario 1 (Member searches and reserves a book):} The member searches for a book, reserves it, and checks it out.
\item
\protect\phantomsection\label{Jinz3BEdi5-y5pU3ZsrwM}
\textbf{Scenario 2 (Member returns the book (on time and late)):} The member returns the book, and the system calculates if a fine is due.
\item
\protect\phantomsection\label{LoLbHdGec5VHdwjFqzwIU}
\textbf{Scenario 3 (Book renewal):} The member renews a borrowed book.
\item
\protect\phantomsection\label{sy_-cBalcSVj99HSEXFEo}
\textbf{Notifications:} Overdue and reservation notifications are simulated.
\end{itemize}
\\
\textbf{Hint: Try customizing the code!}
\\
This code is designed for experimentation, not just reading. You can edit, extend, or modify any part of the example.

\textbf{Executable code: Library management system}

\begin{lstlisting}[language=Java]
import java.util.*;
import java.text.SimpleDateFormat;

public class Driver {
    public static void main(String[] args) {
        try {
            // ============================================
            // >>> SYSTEM INITIALIZATION
            // ============================================
            System.out.println("\n========== SYSTEM INITIALIZATION ==========\n");
            Address libAddress = new Address("1 Main St", "Springfield", "State", 12345, "Country");
            Library library = Library.getInstance("Springfield Public Library", libAddress);

            // ---> Authors
            Author author1 = new Author("Jane Austen", libAddress, "austen@email.com", "123456", "Famous novelist");
            Author author2 = new Author("Mark Twain", libAddress, "twain@email.com", "234567", "Another novelist");

            // ---> Books and BookItems
            SimpleDateFormat sdf = new SimpleDateFormat("yyyy-MM-dd");
            Book book1 = new Book("ISBN123", "Pride and Prejudice", "Novel", "Publisher A",
                    sdf.parse("2000-01-01"), "English", 300, BookFormat.HARDCOVER, Arrays.asList(author1));
            Book book2 = new Book("ISBN456", "Adventures of Huckleberry Finn", "Adventure", "Publisher B",
                    sdf.parse("2005-05-20"), "English", 250, BookFormat.PAPERBACK, Arrays.asList(author2));
            Rack rack1 = new Rack(1, "A1");
            Rack rack2 = new Rack(2, "B2");
            BookItem bookItem1 = new BookItem("BI001", book1, rack1, 30.0, sdf.parse("2020-01-01"), book1.publicationDate);
            BookItem bookItem2 = new BookItem("BI002", book2, rack2, 25.0, sdf.parse("2021-06-15"), book2.publicationDate);
            library.catalog.addBookItem(bookItem1);
            library.catalog.addBookItem(bookItem2);

            // ---> Users
            LibraryCard cardMember = new LibraryCard("CARD1001", new Date());
            Person personMember = new Person("Alice", libAddress, "alice@email.com", "345678");
            Member member = new Member("MEM001", "pass", personMember, cardMember);

            LibraryCard cardLibrarian = new LibraryCard("CARD2001", new Date());
            Person personLibrarian = new Person("Libby", libAddress, "libby@email.com", "456789");
            Librarian librarian = new Librarian("LIB001", "pass", personLibrarian, cardLibrarian);

            // ============================================
            // >>> SCENARIO 1: Search and Reserve
            // ============================================
            System.out.println("\n------------------------------");
            System.out.println(">>> SCENARIO 1: Member searches and reserves a book");
            System.out.println("------------------------------\n");

            List<BookItem> foundBooks = library.catalog.searchByTitle("Pride and Prejudice");
            if (!foundBooks.isEmpty()) {
                BookItem bookItem = foundBooks.get(0);

                System.out.println("-> Member [" + member.person.name + "] attempts to reserve book item [" + bookItem.id + "]");
                member.reserveBookItem(bookItem);

                System.out.println("-> Member [" + member.person.name + "] checks out reserved book item [" + bookItem.id + "]");
                member.checkoutBookItem(bookItem);

            } else {
                System.out.println("Book not found.");
            }

            // ============================================
            // >>> SCENARIO 2: Return Book (Late Return)
            // ============================================
            System.out.println("\n------------------------------");
            System.out.println(">>> SCENARIO 2: Member returns the book late");
            System.out.println("------------------------------\n");

            // Fast forward due date to simulate late return
            Calendar lateCalendar = Calendar.getInstance();
            lateCalendar.setTime(new Date());
            lateCalendar.add(Calendar.DATE, -3); // Simulate due date as 3 days ago
            bookItem1.dueDate = lateCalendar.getTime();

            System.out.println("-> Member [" + member.person.name + "] returns book item [" + bookItem1.id + "] after due date.");
            member.returnBookItem(bookItem1);

            // ============================================
            // >>> SCENARIO 3: Renew Book
            // ============================================
            System.out.println("\n------------------------------");
            System.out.println(">>> SCENARIO 3: Member renews a book");
            System.out.println("------------------------------\n");

            // Check out again for demonstration
            System.out.println("-> Member [" + member.person.name + "] checks out the book again.");
            member.checkoutBookItem(bookItem1);

            System.out.println("-> Member [" + member.person.name + "] renews the book.");
            member.renewBookItem(bookItem1);

            // ============================================
            // >>> NOTIFICATIONS
            // ============================================
            System.out.println("\n------------------------------");
            System.out.println(">>> NOTIFICATIONS");
            System.out.println("------------------------------\n");

            EmailNotification emailNotification = new EmailNotification(
                    "N001", "Your book is overdue!", member.person.email);
            emailNotification.sendNotification();

            PostalNotification postalNotification = new PostalNotification(
                    "N002", "Please return your book!", member.person.address);
            postalNotification.sendNotification();

            System.out.println("\n========== END OF DEMO ==========\n");

        } catch (Exception ex) {
            ex.printStackTrace();
        }
    }
}

\end{lstlisting}

\subsection{Wrapping up}\label{zlfPwLLWjKm88y1gIhvCG}
\\
In this chapter, we've explored the complete design of a library management system. We 've also looked at how a basic library management system can be visualized using various UML diagrams and designed using object-oriented principles and design patterns.

% End of section content